In Corollary~\ref{cor:bounded-operator-spectrum-norm-bound} we used the fact that we could expand $R_\lambda(A)$ in a power series if $|\lambda|>\|A\|$. In fact, we can do much better.

\begin{Theorem}
\label{thm:local-resolvent-power-series}
Let $A\in\mathcal{L}(X)$ and $\lambda_0\in\rho(A)$. Suppose $\lambda\in\mathbb{C}$ lies in the disk
\begin{equation}\label{eq:local-resolvent-series-disk}
|\lambda-\lambda_0|<\dfrac{1}{\|R_{\lambda_0}\|}.
\end{equation}
Then $\lambda\in\rho(A)$ and
\begin{equation}\label{eq:local-resolvent-power-series}
R_\lambda(A)=\sum_{k=0}^\infty(\lambda-\lambda_0)^kR_{\lambda_0}(A)^{k+1}.
\end{equation}

\end{Theorem}

\begingroup
\postdisplaypenalty=10000
\renewcommand{\endProof}{\par\nobreak\hfill$\blacksquare$\par\vspace{0.1cm}}
\begin{Proof}
Let $\lambda_0\in\rho(A)$ and $\lambda\in\mathbb{C}$ satisfying \eqref{eq:local-resolvent-series-disk} be given. We then write\[\begin{aligned}
A-\lambda I&=A-\lambda_0I-(\lambda-\lambda_0)I\\
&=(A-\lambda_0I)[I-(\lambda-\lambda_0)(A-\lambda_0I)^{-1}]\\
&=(A-\lambda_0I)[I-(\lambda-\lambda_0)R_{\lambda_0}(A)],
\end{aligned}\]or simply
\begin{equation}\label{eq:local-resolvent-operator-factorization}
A-\lambda I=(A-\lambda_0I)B,
\end{equation}
where
\begin{equation}\label{eq:local-resolvent-factor-definition}
B:=[I-(\lambda-\lambda_0)R_{\lambda_0}(A)].
\end{equation}
Now since $\|(\lambda-\lambda_0)R_{\lambda_0}(A)\|<1$, we can use Theorem~\ref{thm:neumann-operator-series} to show that $B$ has a bounded inverse and
\begin{equation}\label{eq:local-resolvent-factor-inverse-series}
B^{-1}=\sum_{k=0}^\infty(\lambda-\lambda_0)^kR_{\lambda_0}(A)^k.
\end{equation}
Now we use this and \eqref{eq:local-resolvent-operator-factorization} to get
\begin{equation}\label{eq:local-resolvent-series-factorization-conclusion}
R_\lambda(A)=(A-\lambda I)^{-1}=B^{-1}(A-\lambda_0I)^{-1}=\sum_{k=0}^\infty(\lambda-\lambda_0)^kR_{\lambda_0}(A)^{k+1}.
\end{equation}
This completes the proof.
\end{Proof}
\endgroup

This immediately implies the following.

\begin{Corollary}
\label{cor:bounded-operator-resolvent-set-open}
The resolvent set $\rho(A)\subset\mathbb{C}$ of a bounded operator $A$ is open.
\end{Corollary}

Combining this with Theorem~\ref{cor:bounded-operator-spectrum-norm-bound} and the Heine-Borel theorem gives us another important result.

\begin{Corollary}
\label{cor:bounded-operator-spectrum-compact}
The spectrum $\sigma(A)\subset\mathbb{C}$ of a bounded operator $A$ is a compact set.
\end{Corollary}

We will be able to use the power series representation of Theorem~\ref{thm:local-resolvent-power-series} to employ some elementary techniques of complex variables, but first we need to give a definition of an analytic operator-valued function of a complex variable. The definition we give here holds for a mapping from the complex plane to any Banach space: A mapping to the Banach space of bounded operators $\mathcal{L}(X)$ is a special case.
